\textbf{1. Wave equation.} The wave equation obtained in Problem B.1 is
\[
\boxed{u_{tt}=c^2u_{xx}}\qquad\text{(B.2.1)}
\]
where $u(x,t)$ is the axial displacement and $c$ is the wavespeed.

\textbf{2. Separation of variables.} Assume that the spatial and temporal behaviors can be separated as
\[
u(x,t)=U(x)T(t).
\]

\textbf{3. Differentiation.} The required derivatives are
\[
u_{xx}=U''(x)T(t),
\qquad
u_{tt}=U(x)T''(t).
\]

\textbf{4. Harmonic time dependence.} For vibratory motion with angular frequency $\omega$, use the course-note convention
\[
T(t)=e^{-i\omega t}.
\]

\textbf{5. Harmonic derivatives.} Substitution of this time behavior into the separated solution and its derivatives gives
\[
u(x,t)=U(x)e^{-i\omega t},
\qquad
u_{xx}=U''(x)e^{-i\omega t},
\qquad
u_{tt}=-\omega^2U(x)e^{-i\omega t}.
\]

\textbf{6. Spatial equation before defining the wavenumber.} Substitution into the wave equation gives
\[
c^2U''(x)e^{-i\omega t}=-\omega^2U(x)e^{-i\omega t}.
\]
After canceling the nonzero time factor,
\[
\boxed{c^2U''(x)+\omega^2U(x)=0}\qquad\text{(B.2.6)}
\]

\textbf{7. Wavenumber.} Define the wavenumber by
\[
\boxed{\gamma=\frac{\omega}{c}}\qquad\text{(B.2.7)}
\]

\textbf{8. Spatial ODE.} For $\omega>0$ and $c>0$, dividing the equation from item 6 by $c^2$ gives
\[
\boxed{U''(x)+\gamma^2U(x)=0}\qquad\text{(B.2.8)}
\]

\textbf{9. Exponential trial solution.} Assume
\[
U(x)=e^{s x}\qquad\text{(B.2.9)}
\]

\textbf{10. Characteristic equation.} Then $U''(x)=s^2e^{s x}=s^2U(x)$. Substitution into the spatial ODE gives
\[
\boxed{s^2+\gamma^2=0}\qquad\text{(B.2.10)}
\]

\textbf{11. Characteristic roots.} Solving the characteristic equation gives
\[
\boxed{s_1=+i\gamma,\qquad s_2=-i\gamma}\qquad\text{(B.2.11)}
\]

\textbf{12. Complete spatial solution.} The two independent exponential solutions give the linear combination
\[
\boxed{U(x)=A_1e^{i\gamma x}+A_2e^{-i\gamma x}}\qquad\text{(B.2.12)}
\]
where $A_1$ and $A_2$ are arbitrary complex constants.

\textbf{13. Complete space--time solution.} Restoring the harmonic time factor gives
\[
\boxed{
u(x,t)=A_1e^{i(\gamma x-\omega t)}+A_2e^{-i(\gamma x+\omega t)}
}\qquad\text{(B.2.13)}
\]
The first term has constant phase when $x-ct$ is constant and therefore travels in the positive $x$ direction. The second term has constant phase when $x+ct$ is constant and travels in the negative $x$ direction.

\textbf{14. Trigonometric form and physical displacement.} Euler's formulas give
\[
\begin{aligned}
U(x)
&=A_1(\cos\gamma x+i\sin\gamma x)+A_2(\cos\gamma x-i\sin\gamma x)\\
&=\bigl[i(A_1-A_2)\bigr]\sin\gamma x+(A_1+A_2)\cos\gamma x.
\end{aligned}
\]
Define
\[
C_1=i(A_1-A_2),\qquad C_2=A_1+A_2.
\]
The solution can therefore be written as
\[
\boxed{
U(x)=C_1\sin\gamma x+C_2\cos\gamma x
}
\]
and
\[
\boxed{
u(x,t)=\bigl(C_1\sin\gamma x+C_2\cos\gamma x\bigr)e^{-i\omega t}
}\qquad\text{(B.2.14)}
\]
The measurable displacement is the real part,
\[
\boxed{u_{\mathrm{phys}}(x,t)=\Re\{u(x,t)\}.}
\]
The constants and frequency are selected by initial or boundary data when those data are specified. Problem B.2 considers an infinite tube, so no boundary conditions are imposed and no discrete eigenvalue condition is introduced; the result is the monochromatic wave family parameterized by $\omega$ and the associated $\gamma=\omega/c$.
